\documentclass[11pt,a4paper]{article}
\usepackage[margin=25mm]{geometry}
\usepackage{amsmath,amssymb,amsthm,booktabs}
\newtheorem{lemma}{Lemma}\newtheorem{theorem}{Theorem}
\title{A distinguished maximal subgroup and the order 16 column}
\author{Human Compression project}\date{}
\begin{document}\maketitle
For a finite group $J$, let $x_J(p,q)$ count ordered pairs $(r,s)$
with $|r|=p$, $|s|=q$ and $|rs|=2$, and let $e_J(p,q)$ count those
generating $J$. We first isolate the counting principle used below.

\begin{lemma}\label{distinguished}
Fix a group $G$, an order $q$, and a conjugacy class of maximal
subgroups $\mathcal C$. Suppose every element $s$ of order $q$
lies in exactly one $C(s)\in\mathcal C$. Suppose also that every
proper subgroup generated by a pair of the required type and
not contained in $C(s)$ has a unique maximal overgroup in $G$.
Then
\[
e_G(p,q)=x_G(p,q)-\sum_{M\text{ maximal}}x_M(p,q)
 +\sum_{\substack{M\text{ maximal}\\M\notin\mathcal C}}
 |\{(r,s)\in M^2: |r|=p,\ |s|=q,\ |rs|=2,\ r\in C(s)\}|.
\]
Both sums here range over actual subgroups, not conjugacy classes.
\end{lemma}
\begin{proof}
A nongenerating pair outside its distinguished maximal subgroup
occurs exactly once in the first sum, by hypothesis. A pair inside
that subgroup occurs once there, and once more for each of its
other maximal overgroups. The second sum removes precisely these
additional occurrences. Generating pairs occur in neither sum.
\end{proof}

Now put $G=O'N$ and $q=16$. The primary ATLAS table gives centralizer
order 16 for each of the four element classes of that order. Thus
$C_G(s)=\langle s\rangle$, with unique involution $t=s^8$.
The row 7 maximal subgroups are the involution centralizers.
Exactly one contains $s$, namely $C(s)=C_G(t)$.

\begin{lemma}\label{unique}
Every proper subgroup of $G$ containing an element of order 16
and not contained in a row 7 subgroup has a unique maximal
overgroup in $G$.
\end{lemma}
\begin{proof}
The complete roster has two cyclic order 16 classes, rows 510 and
512. Take their complete ancestor lists and remove all classes
which occur below row 7, as well as $G$ itself. The remaining
sixteen rows are
\[
2,3,4,5,9,10,15,16,21,53,54,56,93,95,99,100.
\]
Rows 2, 3, 15, 16 and 21 are maximal in $G$. Each other row has
exactly one immediate overgroup class, as follows:
\[
\begin{array}{c|rrrrrrrrrrr}
\text{row}&4&5&9&10&53&54&56&93&95&99&100\\\hline
\text{parent}&2&3&4&5&9&10&21&53&54&56&56.
\end{array}                                                     \tag{1}
\]
For each edge $K<L$ in (1), the class lengths and the number of
row $K$ subgroups in a fixed $L$ satisfy
\[
\frac{\#\{L^g:g\in G\}\,
       \#\{\text{row }K\text{ subgroups in }L\}}
      {\#\{K^g:g\in G\}}=1.
\]
This direct incidence calculation proves uniqueness of the actual
parent. Iterating these already proved uniqueness statements
gives the unique maximal overgroup. We do not count arbitrary
paths in the subgroup diagram as distinct containments.

Finally, if the class of a subgroup containing $s$ occurs below
row 7, conjugacy supplies an actual row 7 subgroup containing it.
That subgroup must be $C(s)$ by uniqueness. Thus the class-list
criterion is exactly the condition needed in the statement.
\end{proof}

The only maximal classes containing order 16 elements are
rows 2, 3, 7, 15, 16 and 21. We calculate their contributions
to the correction in Lemma~\ref{distinguished}. For a fixed
maximal subgroup $M$ other than row 7, this correction counts
pairs $(r,s)$ for which $r$ commutes with $s^8$. Equivalently,
it sums $x_{C_M(t)}(p,16)$ over involutions $t$ of $M$ whose
centralizer contains an order 16 element. No pair is repeated:
the sole involution commuting with $s$ is $s^8$.

\begin{lemma}\label{centralizers}
The relevant centralizers, their numbers in a fixed maximal
subgroup, and the numbers of actual maximal subgroups are
\[
\begin{array}{c|c|r|r}
\text{maximal rows}&\text{centralizer}&\text{number in }M&[G:M]\\\hline
2,3&H\ (|H|=1344)&2793&122760\\
15,16&D_{32}&465&30968784\\
21&T\ (|T|=1536)&7&42858585.
\end{array}                                                     \tag{2}
\]
Here $D_{32}$ is dihedral of order 32. The block group $H$ is
the previously accepted inner-involution centralizer in $L_3(7):2$.
The group $T$ is the centralizer of a nonzero involution in the
normal subgroup $C_4^3$ of the accepted row 21 model.
\end{lemma}
\begin{proof}
In a projective maximal, $s^8$ is inner because the quotient by
the socle has order two. The accepted block argument proves
that inner involutions form one class with centralizer $H$.
Their number is $3753792/1344=2793$.

Rows 15 and 16 are $\operatorname{PSL}_2(31)$. An involution lifts
to a determinant-one matrix $A$ with $A^2=-I$. Since $-1$ is
not a square in $\mathbb F_{31}$, this is a nonsplit torus
element. The matrices commuting with it in $\operatorname{SL}_2(31)$
are the norm-one elements of $\mathbb F_{31^2}^\times$, cyclic of
order 32. In the projective quotient they give a cyclic group
of order 16. A projective centralizer may also conjugate $A$ to
$-A$. Such elements exist: Frobenius on the quadratic extension
has determinant $-1$, which is corrected by multiplication with
an element of norm $-1$. The resulting determinant-one matrix
has square $-I$ and acts by inversion on the torus. Thus the
projective centralizer is $D_{32}$. The irreducible characteristic
polynomial makes all these involutions conjugate: a general
linear conjugator can be adjusted to determinant one using the
surjective norm of its field centralizer. Their number is
$14880/32=465$.

For row 21 we reuse the exact permutation model already identified
in the accepted order 12 exclusion. That identification uses the
complete perfect-group catalogue and the acquired character data;
those remain explicit premises, not a new classification claim.
The new literal certificate enumerates its 10752 permutations
from the same retained generators. Its indicated six kernel
generators span a normal abelian subgroup of order 64 with
seven elements of order two and 56 of order four; hence it is
$C_4^3$. The seven nonidentity involutions form one orbit under
the full group's generators. Every one of the 2688 order 16
elements has eighth power in this orbit. Commutation with one
displayed involution gives the complete 1536-element centralizer
$T$, with generators whose closure is exactly that set.
This proves that there are precisely seven relevant centralizers,
all isomorphic, and no other involution type contributes.
\end{proof}

\begin{theorem}
For every source order $p$,
\begin{align*}
e_G(p,16)={}&x_G(p,16)-\sum_M[G:M]x_M(p,16)\\
 &+2\cdot342868680x_H(p,16)\\
 &+2\cdot30968784\cdot465x_{D_{32}}(p,16)
   +300010095x_T(p,16),
\end{align*}
where the sum is over the thirteen maximal conjugacy classes.
\end{theorem}
\begin{proof}
Lemmas~\ref{unique} and~\ref{distinguished} reduce all repeated
maximal contributions to the pairs commuting with $s^8$.
Substitute (2), using $122760\cdot2793=342868680$ and
$42858585\cdot7=300010095$. These are actual parent incidences;
the formula makes no uniqueness assertion about a dihedral
subgroup's global maximal parents.
\end{proof}

The exact finite values relevant to the six previously unpaid
entries are
\[
\begin{array}{c|rrrrrr}
p&3&4&6&8&12&16\\\hline
x_H&2688&2688&8064&10752&5376&6384\\
x_{D_{32}}&0&0&0&0&0&8\\
x_T&6144&6144&6144&12288&12288&8832.
\end{array}
\]
The independent reader reconstructs the full model, kernel and
centralizer from literal permutations, checks all order 16
conjugacy classes, and tests every partner for every representative.
It also reuses the accepted complete block set and constructs the
dihedral group directly on sixteen points. There is no ambient
$O'N$ enumeration or reliance on a software-generated subgroup order.

The acquired exact character evaluations now give
\[
\begin{array}{c|rrrrrrrr}
p&3&4&5&6&7&8&10&11\\\hline
e_G/(2|G|)&68&1292&1966&4700&7448&22040&17822&32488\\[2pt]
p&12&14&15&16&19&20&28&31\\\hline
e_G/(2|G|)&9764&12504&15808&88784&56052&35644&25316&23048.
\end{array}
\]
Division is by $|\operatorname{Aut}(G)|=2|G|$, whose action on
generating ordered pairs is free; there is no duality quotient.
All sixteen entries agree with the fixed printed table. Ten
are cross-checks of earlier accepted proofs and six are newly
recovered here.

\paragraph{Attribution and limits.}
The numerical assertions are from Connor and Leemans,
\emph{Algorithmic enumeration of regular maps}, Table~1.
The complete subgroup roster, named maximal types, character
evaluations and accepted perfect-model identification remain
explicit borrowed premises. The contribution is the distinguished
maximal counting argument, its exact unique-parent application,
and the small centralizer certificate. We make no additional
chirality, priority, open-question, Lean or whole-original claim.
\end{document}
